\chapter{The Avatar System}
\label{ch:avatar}

The Avatar system is the strangest corner of CNA, and deliberately so. Chapter~1 already
flagged it: unlike every other namespace in this book, Avatar was ported from
a \emph{decompiled} real Microsoft XNA~4.0 reference assembly, because FNA never implemented
it at all --- real avatar rendering depended on Xbox Live's cloud avatar-editor service, which
has been offline for well over a decade. This chapter explains what that means concretely:
the base API is not a limitation CNA imposed, but a byte-exact reproduction of behavior that
was \emph{already inert} on the real assembly, and a wholly separate, CNA-original extension
sits alongside it to provide actual, real rendering.

\section{The base API: faithfully, provably inert}

\cnaclass{AvatarRenderer}'s own class documentation states the finding plainly: neither of
its two constructor overloads ever actually reads its arguments, and
\texttt{getStateProperty()} unconditionally forces itself to
\cnaclass{AvatarRendererState::Unavailable} on every single read --- nothing anywhere in the
class ever assigns \texttt{Ready} or \texttt{Loading}. This is not a bug CNA introduced; it is
the real, observed behavior of the actual reference assembly off real Xbox~360 hardware,
where the mesh, texture, and animation data an avatar needs was always streamed at runtime
from a now-dead cloud service and consequently never present at all. \texttt{Draw()} is
correspondingly a permanent no-op in both its overloads, and the \cnaclass{Matrix}-array
overload specifically still enforces a real, meaningful contract despite drawing nothing ---
it hard-throws \texttt{ArgumentException} unless given exactly the real Xbox avatar
skeleton's 71 bone entries. \texttt{getBindPoseProperty()} throws
\texttt{InvalidOperationException} unless \texttt{State == Ready}, which --- per the above ---
is never true in faithful usage, so this call always throws in practice.

\cnaclass{AvatarDescription} tells the same story from a different angle: its constructor
requires exactly 1{,}021 bytes and throws otherwise, a real format that was never publicly
documented and cannot be reverse-engineered from the reference assembly alone.
\texttt{CreateRandom()} --- despite its name --- never actually randomizes anything; both
overloads always return an all-zero, invalid description, and the second overload's
\cnaclass{AvatarBodyType} parameter is not read at all. \texttt{BeginGetFromGamer}/%
\texttt{EndGetFromGamer} complete synchronously (a ``fully synchronous fake-async
operation,'' in the project's own words) and \texttt{EndGetFromGamer} always returns the same
all-zero, invalid description regardless of which gamer was actually passed in.
\cnaclass{AvatarAnimation} completes the pattern: its constructor never reads the requested
\cnaclass{AvatarAnimationPreset} at all, so every instance ends up with the identical 71
zero-valued bone transforms and a zero-length animation clip, no matter which of the 31 named
presets was requested.

None of this is presented apologetically in the project's own documentation, and this book
follows that lead: it is presented as a faithful, verified reproduction of a real, if
surprising, historical API surface --- exactly the kind of specific, checkable claim
Chapter~1 asked this book to preserve rather than round off into a vague ``Avatar
support is limited.''

\section{SkinnedModelEXT: the separate, real-rendering extension}

Real avatar rendering in CNA does not live inside the classes above at all --- it lives in a
structurally independent, entirely \texttt{NOXNA} system, \cnaclass{SkinnedModelEXT}, built
specifically because \cnaclass{Model}/\cnaclass{ModelMesh} (Chapter~12's models-and-meshes
coverage) encode a rigid, per-mesh parent-bone hierarchy --- the wrong shape entirely for
real per-vertex GPU skinning. \cnaclass{SkinnedModelEXT} owns its own bone hierarchy (a
parent-index array, local bind poses, and global inverse bind poses), its own named animation
clips (each a set of per-bone keyframe tracks with linear translation/scale interpolation and
spherical-linear rotation interpolation), and its own mesh parts, deliberately move-only since
it owns real GPU buffer resources. Crucially, its bone count and hierarchy are wholly
independent of the faithful 71-bone \cnaclass{AvatarBone} array above --- a typical
\cnaclass{SkinnedModelEXT} skeleton runs 50 to 65 bones, comfortably inside
\cnaclass{SkinnedEffect}'s (Chapter~13's stock-effects coverage) 72-bone cap, and the
two bone-numbering schemes are never confused or bridged directly.

A small, opt-in surface bridges the two systems without touching the faithful API's own
inert behavior at all: \texttt{AvatarRenderer::EnableRealRenderingEXT(device, model)} builds
an internal \cnaclass{SkinnedEffect}; \texttt{DrawRealEXT(clipName, position, loop)} then
samples the requested clip and issues a real draw through the ordinary
\cnaclass{GraphicsDevice} path, throwing \texttt{InvalidOperationException} if real rendering
was never enabled in the first place. Nothing about this opt-in path changes what
\texttt{Draw()}, \texttt{getStateProperty()}, or any other faithful-API member does --- a game
that never calls \texttt{EnableRealRenderingEXT} sees exactly the inert behavior described
above, unchanged.

\subsection{Appearance: necessarily invented, not reverse-engineered}

\cnaclass{AvatarAppearanceEXT} provides skin/hair/garment tint colors --- and its own
documentation is explicit that this is a CNA invention, not a reconstruction of the real,
proprietary avatar-appearance format, which is permanently inaccessible from the decompiled
assembly alone. It is tint-only in its current form, with no per-part texture or clothing
geometry customization. One default value was itself the subject of a real, confirmed bug:
the default shoe tint was originally near-black, causing avatar shoes to render as a
``featureless pure-black blob'' until fixed through direct pixel-sampling verification rather
than a guess.

\subsection{Real bugs found building the content pipeline}

A dedicated \texttt{.skinnedmodel.json}/\texttt{.skeleton.bin}/\texttt{.clip.bin} content
pipeline loads \cnaclass{SkinnedModelEXT} assets, and building it end to end surfaced three
genuine, worth-naming bugs. Manifest asset paths originally resolved against the content
root rather than the manifest file's own directory, breaking any subdirectory-organized asset
bundle --- fixed to resolve relative to the manifest itself. A genuine C\texttt{++}
unspecified-evaluation-order bug corrupted keyframe data: chaining several
\texttt{reader.Read<float>()} calls directly as constructor arguments for a
\cnaclass{Vector3}/\cnaclass{Quaternion} scrambled component order, because C\texttt{++} does
not guarantee left-to-right evaluation of function arguments the way a reader might assume ---
fixed by reading each component into a named local first. And the asset-conversion tooling's
own bone-hierarchy reordering step needed a corresponding remap of inverse bind matrices and
per-vertex joint indices, which were still expressed in glTF's own original joint order rather
than the tool's newly reordered one.

\section{Backend support and real content proof}

Real GPU skinning through this extension is proven end-to-end via pixel readback on EasyGL;
Vulkan and BGFX both have the real pipeline wiring in place (descriptor sets and per-frame
bone uniform buffers) but had not, as of the source material this book draws from, been
separately smoke-tested for this specific feature. On \texttt{SDL\_RENDERER}, any attempt to
create the 3D resources this extension needs throws the same ``does not support 3D'' error
already covered in Part~IV's backend chapters, consistent with that backend's documented
scope. Beyond synthetic tests, a dedicated demo application loads real content and renders a
complete, correctly proportioned, animated humanoid avatar in a real window, toggling between
clips live --- alongside seven sibling demo applications covering an animation gallery,
appearance-tint studio, wardrobe hot-swapping, multi-attach stress testing, and a
side-by-side comparison against the faithful, inert base API, so a reader can see both halves
of this chapter's story running at once.

\section{The shape of this chapter, restated}

\begin{center}
\begin{tabular}{lll}
\toprule
& \textbf{Base Avatar API (faithful)} & \textbf{SkinnedModelEXT (\texttt{NOXNA})} \\
\midrule
Purpose & Byte-exact port of real XNA~4.0 & CNA-original, opt-in real renderer \\
Bone rig & Fixed 71-bone array, always zero & Independent count (\raisebox{0.5ex}{\texttildelow}50--65), real skeleton \\
Draw() & Permanent no-op & Real \cnaclass{GraphicsDevice}/\cnaclass{SkinnedEffect} draw \\
Enabled by default & Yes & No --- requires explicit opt-in \\
\bottomrule
\end{tabular}
\end{center}
